\section{Isolation, distinct sums, and prime collisions}
\label{sec:setup}

We first reduce detection to finding an isolated exact sum.  We then give
an exact preliminary search that succeeds when a random half of the
coordinates has sufficiently few distinct subset sums.  The complementary
case supplies the combinatorial information needed later to recover
modular aliases.  Finally, we record a random-prime collision bound
for the nonzero integer discrepancies arising in alias extraction.

All logarithms are to base~$2$, except for $\ln$.  Write
$[n]=\{1,\ldots,n\}$, and for weights $v_1,\ldots,v_n$ write
$v(I)=\sum_{i\in I}v_i$ for $I\subseteq[n]$.  Empty sums are zero.
All constants and decimal parameters in the proof are fixed, independently
of the input.

\subsection{Isolating one exact target}

The following application of the Isolation Lemma of
Mulmuley, Vazirani, and Vazirani~\cite{MVV1987} allows a complex-valued
checksum to detect a solution without an anti-concentration argument.
Choose a power of two $N_0$ with $12n\le N_0<24n$, and draw independent
uniform integers $u_i\in\{1,\ldots,N_0\}$.  Define
\begin{equation}
\label{eq:transformed}
 C=1+nN_0,\qquad c_i=Ca_i+u_i,\qquad
 t_s=Ct+s\quad(0\le s\le nN_0).
\end{equation}
The transformed weights and targets have $O(b+\log n)$ bits, and there
are $nN_0+1=O(n^2)$ targets.

\begin{lemma}[Isolation]
\label{lem:isolation}
For every choice of the tie weights, there is a subset $I\subseteq[n]$
with $a(I)=t$ if and only if there are $s\in\{0,\ldots,nN_0\}$ and
$I\subseteq[n]$ with $c(I)=t_s$.  If the original instance has a
solution, then with probability at least $11/12$ at least one transformed
target $t_s$ has exactly one representing subset.
\end{lemma}

\begin{proof}
The identity $c(I)=t_s$ gives
$C(a(I)-t)=s-u(I)$.  Since $s,u(I)\in[0,nN_0]$ and $C>nN_0$, this
forces $a(I)=t$ and $u(I)=s$.  Conversely, every original solution
represents the transformed target indexed by $s=u(I)$.

Let $\mathcal F=\{I\subseteq[n]:a(I)=t\}$, and suppose
$\mathcal F\ne\varnothing$.  Fix an index $i$ and all tie weights other
than $u_i$.  The minimum tie weight among members of $\mathcal F$
containing $i$ has the form $A_i+u_i$, while the minimum among members
not containing $i$ is a constant $B_i$.  A missing family can be assigned
minimum $+\infty$.  If both families are nonempty, their two minima
coincide for at most one value of $u_i$.  Consequently, the probability
that they coincide is at most $1/N_0$.

If $u$ has two distinct minimum-weight members of $\mathcal F$, they
differ at some index $i$, and the two minima just described coincide
there.  A union bound shows that the minimum is unique with probability
at least $1-n/N_0\ge11/12$.  Its tie weight $s$ then indexes an exact
bin containing just this subset, by the first part of the proof.
\end{proof}

\subsection{Deficiency and the preliminary search}

For the rest of this section, fix the transformed weights $c_i$.
For $V\subseteq[n]$ define its subset-sum support and its deficiency by
\begin{equation}
\label{eq:deficiency}
 \Sigma(V)=\{c(I):I\subseteq V\},\qquad
 d(V)=|V|-\log|\Sigma(V)|.
\end{equation}
Thus $d(V)$ is the exponent saved by storing distinct subset sums
instead of all $2^{|V|}$ subsets: $|\Sigma(V)|=2^{|V|-d(V)}$.
Large deficiency means that many subsets have the same sum.
These definitions make sense for arbitrary integer weights, not only
the positive transformed weights.

\begin{lemma}[Properties of deficiency]
\label{lem:deficiency}
For arbitrary fixed integer weights, $d(\varnothing)=0$ and $d(V)\ge0$.
For $i\notin V$,
\[
 0\le d(V\cup\{i\})-d(V)\le1.
\]
For disjoint $U,V\subseteq[n]$,
\[
 d(U\cup V)\ge d(U)+d(V).
\]
If $V,V'$ have equal size and differ by exchanging one index, then
$|d(V)-d(V')|\le1$.
\end{lemma}

\begin{proof}
There are at most $2^{|V|}$ subset sums, and the empty sum always
belongs to the support.  Adding $i$ replaces $\Sigma(V)$ by
$\Sigma(V)\cup(c_i+\Sigma(V))$, whose cardinality lies between
$|\Sigma(V)|$ and $2|\Sigma(V)|$.  This proves the increment bounds.
For disjoint $U,V$,
$\Sigma(U\cup V)=\Sigma(U)+\Sigma(V)$ has cardinality at most
$|\Sigma(U)|\,|\Sigma(V)|$, giving superadditivity.  Finally, in an
exchange write $V=U\cup\{i\}$ and $V'=U\cup\{j\}$.  Both
deficiencies lie in $[d(U),d(U)+1]$.
\end{proof}

\begin{lemma}[Random restrictions]
\label{lem:half-variance}
Let $V_q$ be a uniformly random $q$-element subset of $[n]$, for
$0\le q\le n$.  Then
\[
 \operatorname{Var}d(V_q)\le q.
\]
Moreover, $\mathbb E d(V_q)$ is nondecreasing in $q$.
\end{lemma}

\begin{proof}
Expose an ordered sample $X_1,\ldots,X_q$ without replacement, and let
\[
 M_r=\mathbb E\bigl[d(\{X_1,\ldots,X_q\})
                   \mid X_1,\ldots,X_r\bigr]
 \quad(0\le r\le q).
\]
Fix a possible prefix of length $r-1$ and compare two possible choices
$i,j$ for $X_r$.  Couple the remaining ordered samples by interchanging
$i$ and $j$.  The resulting final sets either coincide, if the remainder
contains the other index, or differ by the exchange $i\leftrightarrow j$.
By Lemma~\ref{lem:deficiency}, their deficiencies differ by at most~$1$.
Thus the possible conditional values of $M_r$, given the prefix of
length $r-1$, lie in an interval of length at most~$1$.  Since $M_{r-1}$
is their mean, $|M_r-M_{r-1}|\le1$.  The martingale differences are
orthogonal in $L^2$, so
\[
 \operatorname{Var}d(V_q)
 =\sum_{r=1}^q\mathbb E(M_r-M_{r-1})^2\le q.
\]
The case $q=0$ is immediate.  To compare expectations for consecutive
sizes, use the first $q$ and first $q+1$ entries of one uniform random
permutation, and apply the monotonicity in Lemma~\ref{lem:deficiency}.
\end{proof}

We now use distinct sums in the meet-in-the-middle algorithm of
Horowitz and Sahni~\cite{HorowitzSahni1974}.  Related distinctions between
small subset-sum supports and the complementary regime appear in earlier
Subset Sum algorithms~\cite{AKKN2016,NederlofWegrzycki2021}.
The procedure can stop without deciding the instance; that outcome
means only that the later algorithm must be run.

\begin{proposition}[Capped preliminary search]
\label{prop:preliminary}
Given fixed transformed weights and targets from
\eqref{eq:transformed}, there is a randomized procedure that either
returns the correct decision for the original instance or returns
\textup{CONTINUE}.  On every execution its running time is at most
\[
 2^{0.489985n}\operatorname{poly}(n+b+2).
\]
If $V$ is uniform among the $\lfloor n/2\rfloor$-element subsets and
$\mathbb E d(V)>0.01004n$, then the probability of \textup{CONTINUE}
is $O(1/n)$, uniformly in the fixed weights.
\end{proposition}

\begin{proof}
Choose a uniform $\lfloor n/2\rfloor$-element set $V$, returning
\textup{CONTINUE} if the capped sampler described below fails.
Construct the two sorted lists of distinct subset sums on
$V,[n]\setminus V$, one side at a time.  Starting with $\{0\}$, insert an
index $i$ by merging the current list with its translate by $c_i$ and
removing duplicate values.  Cap every list at
\[
 B_n=2^{\lceil0.489985n\rceil}.
\]
If a merge produces more than $B_n$ distinct entries, stop immediately
and return \textup{CONTINUE}.  Otherwise, for every $s$ check the two
final sorted lists for a pair summing to $t_s$.  A linear scan suffices
for each target.  Return \textup{YES} exactly when such a pair exists.
Lemma~\ref{lem:isolation} proves correctness whenever this search
finishes.  Each merge takes $O(B_n)$ list operations before completion
or interruption, and there are only polynomially many merges and
target scans.  The bit lengths noted after \eqref{eq:transformed}
give the stated running time.  Intermediate support sizes never exceed
the final support size, so a side completes whenever its final support
has at most $B_n$ entries.

To sample $V$, choose its elements successively without replacement.
When $k$ indices remain, use the low $\lceil\log k\rceil$ bits of a
fresh random word to propose an index, rejecting values at least $k$.
Allow at most $n^2$ proposals per choice and return \textup{CONTINUE}
if all are rejected.  The sampler uses at most $n^3$ proposals and its
total failure probability is at most $n2^{-n^2}$.  Conditional on no
failure, $V$ is uniform because the success probability of each choice
depends only on $k$, not on the selected index.

Suppose now that $\mathbb E d(V)>0.01004n$.  Both side sizes have
expected deficiency greater than $0.01004n$, by
Lemma~\ref{lem:half-variance}; each side has variance at most $n$.
Chebyshev's inequality therefore gives, for either side $U$,
\[
 \Pr\{d(U)<0.010025n\}
 \le \frac{n}{(0.000015n)^2}=O(1/n).
\]
When neither exceptional event occurs,
\[
 \log|\Sigma(U)|=|U|-d(U)
 \le 0.489975n+\tfrac12
 \le 0.489985n
\]
for all sufficiently large $n$.  Both lists then fit under the cap,
and the procedure decides the instance.  A union bound, including the
sampler-failure probability $n2^{-n^2}$, proves the assertion.
\end{proof}

The expectation in Proposition~\ref{prop:preliminary} is used only in
the analysis; the algorithm does not evaluate it.  For each fixed
choice of the tie weights, either the preliminary search decides with
probability $1-O(1/n)$, or
\begin{equation}
\label{eq:low-deficiency}
 \mathbb E_{\substack{V\subseteq[n]\text{ uniform}\\
                      |V|=\lfloor n/2\rfloor}}d(V)
 \le0.01004n.
\end{equation}
The alias-recovery argument will use this latter inequality.
A \textup{CONTINUE} outcome is not a certificate that it holds.

\subsection{A collision bound for random primes}

Reduction modulo a random prime makes any fixed false integer equality
unlikely.  The following elementary bound is the only number-theoretic
estimate needed for that purpose.

\begin{lemma}[Random-prime collisions]
\label{lem:prime-collision}
There are absolute constants $K$ and $L_*$ such that the following
holds.  For an integer $L\ge L_*$, let $p$ be uniform among the primes
in $[2^L,2^{L+1})$.  If $D$ is a nonzero integer with $|D|<2^B$, where
$B\ge1$, then
\[
 \Pr\{D\equiv0\pmod p\}\le KBL\,2^{-L}.
\]
The same bound applies to the collision of two fixed unequal integers
whose difference has absolute value less than $2^B$.
\end{lemma}

\begin{proof}
The prime number theorem gives at least $c_0 2^L/L$ primes in the
specified interval for an absolute $c_0>0$ and all sufficiently large
$L$.  An explicit bound of this form is also given by
Rosser and Schoenfeld~\cite[Corollary~3]{RosserSchoenfeld1962}.
The nonzero integer $D$ has fewer than $B$ distinct prime divisors,
since their product is at most $|D|$ and each is at least~$2$.
Dividing this count by the number of available primes proves the
claim.  For two unequal integers, apply it to their difference.
\end{proof}
